% ============================================================================
\clearpage
\section*{v4.9.0-r20 Source-Only: A Provenance Hypothesis for \texorpdfstring{\(\Delta_\gamma<1/5\)}{Delta gamma < 1/5}, the Photon-Frame Clock Gain at the Hopf Point}
\addcontentsline{toc}{section}{v4.9.0-r20 Source-Only: A Provenance Hypothesis for Delta gamma < 1/5, the Photon-Frame Clock Gain at the Hopf Point}
\label{sec:v490-r20-delta-gamma-provenance}
% ============================================================================

\paragraph{Purpose.}
r12 derived \(\Delta_\gamma<1/5\) from the chart data of r14 under the premises (P1)--(P4).  It left open where those
data, and therefore the five, come from.  The r65 text is lost.  On 2026-10-22 the author suggested that the value was produced
in a photon-frame calculation, and that it need not be correct or generally valid.  This layer proposes a provenance on that
basis and examines how far the value holds away from it.  The script is \texttt{llm/k59\_herkunft\_delta\_gamma.py}.  The
proposal is a hypothesis about origin.  It does not recover r65, makes no new prediction, and compares nothing with data.

\noindent\textbf{Source-only status.}
This v4.9.0-r20 layer is source-only.  It does not build a PDF, does not assign a new DOI,
and does not authorize a public release.  The v4.9.0 Version DOI
\href{https://doi.org/10.5281/zenodo.23534954}{10.5281/zenodo.23534954} and the Concept DOI
\href{https://doi.org/10.5281/zenodo.19210728}{10.5281/zenodo.19210728} are unchanged by it.

\subsection*{v4.9.0-r20 the hypothesis}
\label{subsec:v490-r20-hypothesis}

r14 calls the charts with a window ``photon-frame timing-assisted'' (Proposition~\ref{prop:v470-r14-finite-flowrate-classification}).
The gain \(G\) is therefore the gain of the photon-frame clock.  A clock frequency converts between frames with the frame factor
(Corollary~\ref{cor:gue-frame}).

\paragraph{Theoremlet v4.9.0-r20.1.}
At the Hopf point \(\bsym^2=1/3\), the doublet state of r17 has \(\gL=2\), \(p_+=3/4\) and \(p_-=1/4\).  The binding chart \(V3\) of
r12 and the clock seed of r15 consist of exactly these numbers:
\[
G_{V3}=2=\gL,\qquad \kappa_{V3}=\tfrac14=p_-,\qquad
N_{V3}=5=1+\frac1{p_-}=2+\frac1{\bsym^2},\qquad \kappa_{\rm seed}=\tfrac34=p_+ .
\]
The value \(N=5\) is the balance condition \(\kappa(N-1)=1\) of r12 at \(\kappa=p_-\).  With these values
\(1-16/(GS)=1-16/(2\cdot10)=1/5\) exactly.  The collar row of r05 (ratio \(5/6=1/(1+\tfrac15)\)) and the C5 factor \(3/4=p_+\)
(r17) belong to the same point.

\paragraph{Proof sketch.}
r17.1 with \(\bsym^2=1/3\): \(\gL=(1+\bsym^2)/(1-\bsym^2)=2\), \(p_-=\bsym^2/(1+\bsym^2)=1/4\).  Exact rational arithmetic.  \(\square\)

\paragraph{Hypothesis.}
The bound \(\Delta_\gamma<1/5\) was obtained with the photon-frame clock gain set to the frame factor at the Hopf point,
\(\gL=2\), the continuum value \(\bsym=1/\sqrt3\) without the lattice correction of Remark~\ref{conj:cl-correction}.  A rounding
of \(\gL=1.864\) to the half-integer grid of r14 gives the same \(G=2\).  The hypothesis explains why the five arises at
exactly one chart (r92).  It also explains why the C5 factor and the clock bound share one point.

\paragraph{Weight of the evidence.}
The r14 values lie on grids of quarters (\(\kappa\)) and halves (\(G\)).  A single coincidence such as \(G=2\) would carry little
weight.  Four numbers that match at one point, together with the balance condition, carry more.  Against the hypothesis:
the other window charts \(V5\) (\(G=5/2\)) and \(V8\) (\(G=2\), \(\kappa=1\)) have no comparable reading, and the threshold \(16\) is
dyadic and frame-free.

\subsection*{v4.9.0-r20 the bound away from the Hopf point}
\label{subsec:v490-r20-away}

\paragraph{Theoremlet v4.9.0-r20.2.}
Keep the chart form of r12 and the threshold \(16\), and replace the Hopf values by those at the measured offset.  Then the bound
of \(V3\) becomes:
\begin{center}
\small
\begin{tabular}{@{}>{\raggedright\arraybackslash}p{0.50\textwidth}rrr@{}}
\toprule
reading & \(G\) & \(\kappa\), \(N\) & bound\\
\midrule
Hopf point (source) & \(2\) & \(1/4\), \(5\) & \(0.2000\)\\
\(G=\gL\), chart fixed & \(1.8644\) & \(1/4\), \(5\) & \(0.1418\)\\
\(G=e^\eta\) (clock as scale quantity, r19), chart fixed & \(1.8542\) & \(1/4\), \(5\) & \(0.1371\)\\
\(G=\gL\), \(\kappa=p_-=0.2318\), \(N=5\) & \(1.8644\) & \(0.2318\), \(5\) & \(0.1095\)\\
\(G=e^\eta\), \(\kappa=p_-\), \(N=5\) & \(1.8542\) & \(0.2318\), \(5\) & \(0.1045\)\\
\(G=\gL\), balanced \(N=1+1/p_-=5.31\) & \(1.8644\) & \(0.2318\), \(5.31\) & \(0.1925\)\\
\bottomrule
\end{tabular}
\end{center}
In every reading the bound lies below \(1/5\).  A different generalization, \(\tan\alpha<1/\gL\) in place of \(\tan\alpha<1/2\),
gives \(1/(1+\gL^2)=0.2234\) instead and lies above it.

\paragraph{Proof sketch.}
\(\Delta_\gamma^{\max}=1-16/(GS)\) with \(S=N+\kappa N(N-1)\) and the values of r17.1 at \(\bsym=\tanh\eta\).  Numerical.  \(\square\)

\paragraph{Consequences.}
\begin{enumerate}[leftmargin=2em]
\item \(1/5\) is a value of the Hopf point.  It is not frame-invariant.  At the measured offset the bound depends on how the
chart is carried over, and the readings that keep the chart form give \(0.10\) to \(0.19\).
\item In those readings \(1/5\) is not conservative.  A drift loss between the true bound and \(1/5\) would be admitted by the
stated condition and still close the \(V3\) window.
\item r12.2(b) changes with it.  For the readings with the chart fixed, the bound is about \(0.14\).  The tolerated collar row
of r05, which exports \(1/6=0.167\), then lies outside, and so does the C5 stage (\(1/4\)).  The separation of the C5 stage from
the clock selector stays, but the reading that r05 tolerated that row for this reason holds only at the Hopf point.
\end{enumerate}

\subsection*{v4.9.0-r20 guard status}
\label{subsec:v490-r20-guard}

\(W_{\rm clk}^{\rm loc}\) under \(\Delta_\gamma<1/5\) is a permitted local statement.  This layer makes its bound frame-typed:
\(\Delta_\gamma<\Delta_\gamma^{\max}(\eta)\), with \(\Delta_\gamma^{\max}=1/5\) at the Hopf point.  At the measured offset the
condition becomes narrower, about \(0.14\) for the chart fixed.  A narrower local condition claims less, not more.  No guard is
spent, and nothing global is approached.  Which carrying-over is right is not decided here.  It needs the clock dynamics that
r13 left open.

\subsection*{v4.9.0-r20 register entries}
\label{subsec:v490-r20-register}

\begin{center}
\small
\begin{tabular}{@{}>{\raggedright\arraybackslash}p{0.34\textwidth}>{\raggedright\arraybackslash}p{0.58\textwidth}@{}}
\toprule
statement & status after v4.9.0-r20\\
\midrule
provenance of \(\Delta_\gamma<1/5\) & hypothesis: photon-frame clock gain \(G=\gL=2\) at the Hopf point; \(V3\) and clock seed are Hopf data (r20.1); r65 text still lost\\
\(\Delta_\gamma<1/5\) & exact at the Hopf point; frame-typed, at the measured offset \(0.10\)--\(0.19\) with the chart form kept (r20.2)\\
r12.2(b), collar row of r05 & inside only at the Hopf point; outside at about \(0.14\)\\
Hopf point & now carries the C5 factor \(3/4\), the export \(1/4\), the \(V3\) coupling and gain, and \(\Delta_\gamma=1/5\)\\
\bottomrule
\end{tabular}
\end{center}

\subsection*{v4.9.0-r20 ledger and guard update}
\label{subsec:v490-r20-ledger-guard}

\[
R_{\rm prov}^{v490\text{-}r20}=(0,0,1)
\]
for (provenance hypothesis stated with its evidence, bound away from the Hopf point computed, choice of carrying-over open).

The layer spends no guard:
\[
R_{\rm glob}^{(v490\text{-}r20,+)}=(0,1,1,1,1),
\qquad
\Gamma_{\rm guard}^{v490\text{-}r20}=(1,1,1,1,1),
\qquad
\beta_4=1/256\ \text{unspent}.
\]
It does not imply exact finite \(C5\) zero closure, exact all-readout terminality,
\(C6/T2\), a global photon-clock/RH theorem, or global all-readout closure.

